% Generated draft: requires mathematical and source review.
\documentclass[11pt]{article}
\usepackage[margin=25mm]{geometry}
\usepackage{fontspec}
\setmainfont{Latin Modern Roman}
\usepackage{amsmath,amssymb,mathtools,unicode-math}
\setmathfont{Latin Modern Math}
\usepackage{enumitem,needspace,xurl,hyperref}
\setlength{\parindent}{0pt}
\setlength{\parskip}{5pt}
\setlength{\emergencystretch}{3em}
\newcommand{\N}{\mathbb{N}}
\newcommand{\Z}{\mathbb{Z}}
\newcommand{\Q}{\mathbb{Q}}
\newcommand{\R}{\mathbb{R}}
\newcommand{\C}{\mathbb{C}}
\newcommand{\F}{\mathbb{F}}
\newcommand{\A}{\mathbb{A}}
\newcommand{\PP}{\mathbb{P}}
\newcommand{\E}{\mathbb{E}}
\newcommand{\eps}{\varepsilon}
\newcommand{\dd}{\,\mathrm{d}}
\begin{document}
\section*{UnsolvedMath ID: 20000220}
\textbf{Source problem code:} catalogue entry
\subsection{Definitions and mathematical statement}
The supplied verification says the preceding workshop item defines \(X\subset\PP^2\) as an elliptic curve, hence a smooth plane cubic. For an integer \(d\geq1\), the original asks for the number of degree-\(d\) rational plane curves meeting \(X\) at exactly one point. For an irreducible curve \(C\ne X\), the literal condition is \(\operatorname{Supp}(C\cap X)=\{p\}\). This does not by itself require a single branch over \(p\) on the normalization. A one-place maximal-contact condition instead requires the normalization map \(\nu:\PP^1\to C\) to satisfy \(\nu^*X=3d\,q\) for one point \(q\), with \(\nu(q)=p\). The report warns that these counting conventions, logarithmic stable-map lengths, and counts of distinct reduced image curves must be separated.

In the characteristic-zero elliptic-curve interpretation, choose a flex origin \(o\in X\), so \(\mathcal O_X(1)\cong\mathcal O_X(3o)\). A maximal contact point must satisfy \(3d(p-o)=0\). Define its minimal index to be the smallest \(k\geq1\) with \(3k(p-o)=0\). For such points in degree \(d\), \(k\mid d\). The report's candidate torsion audit gives
\[
A_k=9J_2(k),\qquad
J_2(k)=k^2\prod_{\ell\mid k,\ \ell\ \mathrm{prime}}(1-\ell^{-2}).
\]
These count available contact points of each minimal index, not automatically distinct curves.

The report attributes to Bousseau's all-degree theorem the fixed primitive degree-nine log-BPS invariant, equivalently the length of the finite primitive log-map moduli scheme,
\[
\Omega_9=1{,}027{,}737.
\]
Primitive here means minimal contact index \(9\). This value does not by itself count distinct reduced image curves, since scheme length need not equal the number of reduced points. Using the report's normalization for the fixed-primitive invariants \(\Omega_e\) and total unspecified-contact relative invariant \(N_d\), its aggregate audit identity is
\[
(-1)^{d-1}N_d=\frac9{d^2}\sum_{e\mid d}e^4\Omega_e.
\]
The exact virtual-count and sign normalizations defining these symbols are not otherwise supplied.

In degree nine the possible minimal indices and their contact-point counts are
\[
\begin{array}{c|ccc}
k&1&3&9\\\hline
A_k&9&72&648
\end{array}
\]
and the asserted all-primitive aggregate length and full unspecified-contact relative invariant are respectively
\[
648\cdot1{,}027{,}737=665{,}973{,}576,\qquad
N_9=749{,}220{,}300+\frac19.
\]
The fractional total illustrates why it cannot simply be the literal number of distinct image curves. The report's three-stratum distinction does not remove the original branch, reducedness, and counting ambiguities.

The workshop text says the answer was known for \(d\leq8\) ``by Takashi,'' that \(d=9\) would be new and interesting, and that the motivation is local mirror symmetry; it suggests degenerating \(X\) to three lines and asks about a tropical approach. These are retained as historical workshop remarks, not a current literature assessment.
\subsection{Short English statement}
Count rational plane curves meeting a smooth cubic at one point, while distinguishing literal image curves from one-place logarithmic counts. The report gives a degree-nine primitive invariant and contact-stratum audit, but does not identify these lengths with counts of distinct reduced curves.
\subsection{Sources}
\begin{itemize}
\item AIM workshop ``Degenerations in algebraic geometry,'' Problem 3.1: \url{https://aimath.org/pastworkshops/degenalggeomproblems.pdf}. The supplied verification recovers the smooth plane cubic \(X\) from the preceding item; the PDF was not accessed here.
\item Local canonical registry \texttt{aim-algebraic-geometry-notes.json}, \texttt{notes[219]}, tag \texttt{problem}; permanent UnsolvedMath ID \(20000220\), source code \texttt{catalogue entry}. Original gives the one-point question, the historical \(d\leq8\) attribution ``Takashi,'' and the degeneration/tropical suggestions. The extracted trailing ``4'' is an apparent pagination artifact, not a degree condition. Original remarks and literature field are otherwise empty.
\item Supplied research attempt 1 attributes the earlier count to Takahashi and the all-degree logarithmic result to Bousseau. No titles, identifiers, or full authoritative references are supplied. The torsion audit is a candidate in this machine-generated, non-peer-reviewed aid requiring independent verification before citation.
\end{itemize}
\end{document}
